% Generated draft: requires mathematical and source review.
\documentclass[11pt]{article}
\usepackage[margin=25mm]{geometry}
\usepackage{fontspec}
\setmainfont{Latin Modern Roman}
\usepackage{amsmath,amssymb,mathtools,unicode-math}
\setmathfont{Latin Modern Math}
\usepackage{enumitem,needspace,xurl,hyperref}
\setlength{\parindent}{0pt}
\setlength{\parskip}{5pt}
\setlength{\emergencystretch}{3em}
\newcommand{\N}{\mathbb{N}}
\newcommand{\Z}{\mathbb{Z}}
\newcommand{\Q}{\mathbb{Q}}
\newcommand{\R}{\mathbb{R}}
\newcommand{\C}{\mathbb{C}}
\newcommand{\F}{\mathbb{F}}
\newcommand{\A}{\mathbb{A}}
\newcommand{\PP}{\mathbb{P}}
\newcommand{\E}{\mathbb{E}}
\newcommand{\eps}{\varepsilon}
\newcommand{\dd}{\,\mathrm{d}}
\begin{document}
\section*{UnsolvedMath ID: 11100046}
\textbf{Source problem code:} catalogue entry
\subsection{Definitions and mathematical statement}
In the standard finite-group setting, a Mackey functor for \(G\) assigns abelian groups to finite \(G\)-sets with compatible restriction, transfer, additivity and double-coset identities. A Green functor \(R\) is a ring object for the Mackey box product, and an \(R\)-module is a Mackey functor with compatible \(R\)-action. The first task is to construct an appropriate derived category of these modules and analyze it as an algebraic model for equivariant stable homotopy. A natural candidate is the localization of chain complexes of \(R\)-modules at homology isomorphisms, but boundedness, group class and the exact derived/model structure are not prescribed by the source.

The requested analysis particularly asks for analogues of derived-Noetherian-ring classification results for:
\begin{itemize}
\item thick subcategories: full triangulated subcategories closed under retracts;
\item localizing subcategories: triangulated subcategories closed under coproducts;
\item colocalizing subcategories: triangulated subcategories closed under products.
\end{itemize}
Their object scope and any tensor-ideal condition must be specified. In the familiar commutative-Noetherian-ring analogy, thick classification usually concerns compact objects, not arbitrary thick subcategories of the unbounded derived category; this qualification is not repeated in the extract. General Green functors need not be commutative or satisfy the needed Noetherian/support hypotheses, so no universal spectrum classification is asserted without them.

The source points to Neeman and Hovey--Palmieri--Strickland as the model and says Gaunce Lewis had considered the problem. The summary reports that derived Mackey-module categories have since been constructed and partially analyzed, while distinguishing this from a full classification.
\subsection{Short English statement}
Construct and study the derived category of modules over a Green functor, seeking correctly scoped analogues of thick, localizing and colocalizing subcategory classifications.
\subsection{Sources}
\begin{itemize}
\item Mark Hovey, \emph{Algebraic Topology Problem List}, ``Equivariant homotopy,'' item 4, surviving mirror \url{https://www-users.cse.umn.edu/~tlawson/hovey/}. Recorded access 2026-07-29; extraction html-ordered-list.
\item Amnon Neeman's source-named derived-category paper; Hovey--Palmieri--Strickland for alternative proofs; Gaunce Lewis, credited with considering the problem. Complete bibliography and exact theorem hypotheses are absent.
\item Local UnsolvedMath registry, \texttt{data/unsolvedmath/problems.json}, permanent ID \(11100046\), exact source code \texttt{catalogue entry}; supplied partially\_solved/PARTIAL-PROGRESS summary reports construction progress but partial classification.
\end{itemize}
\end{document}
